% Section declaración - SECTOR PÚBLICO / ENTIDADES GUBERNAMENTALES
\par{
	
	Economista comprometido con el servicio público y el desarrollo del país. Con formación sólida en economía, análisis de datos y gestión pública, busco contribuir al fortalecimiento institucional y a la toma de decisiones basadas en evidencia para mejorar la calidad de vida de los ciudadanos.
	
	Experiencia en elaboración de informes técnicos, procesamiento de bases de datos del sector público, análisis de indicadores sociales y económicos, y seguimiento de políticas públicas. Familiarizado con sistemas de información estatal (SIAF, SIGA, SEACE) y marco normativo del sector público peruano.
	
	Busco la oportunidad de formar parte de una institución pública comprometida con la transparencia, eficiencia y servicio a la ciudadanía, donde pueda aportar mis conocimientos técnicos, capacidad analítica y valores de integridad para contribuir al logro de objetivos institucionales y al bienestar común.

	%%%%%%%%%%%%%%%%%%%%%%%%%%%%%%%%%%%%%%%%
	% VERSIÓN ALTERNATIVA - Análisis de políticas públicas
	%%%%%%%%%%%%%%%%%%%%%%%%%%%%%%%%%%%%%%%%
	%Economista especializado en análisis y evaluación de políticas públicas con enfoque en desarrollo social y económico. Experto en diseño de indicadores, seguimiento de metas gubernamentales, evaluación de impacto y elaboración de documentos técnicos para gestión pública basada en resultados.
	
	%Dominio de metodologías de evaluación ex-ante y ex-post, análisis costo-beneficio, teoría del cambio y marco lógico. Experiencia en trabajo con microdatos de programas sociales, análisis de focalización y estudios de línea de base. Capacidad de redactar términos de referencia, informes de evaluación y notas de política pública.
	
	%Busco contribuir en áreas de planificación estratégica, presupuesto por resultados o unidades de evaluación donde pueda aplicar herramientas de economía aplicada para fortalecer el ciclo de políticas públicas y mejorar la efectividad del gasto público en beneficio de la población.

	%%%%%%%%%%%%%%%%%%%%%%%%%%%%%%%%%%%%%%%%
	% VERSIÓN ALTERNATIVA - Desarrollo regional
	%%%%%%%%%%%%%%%%%%%%%%%%%%%%%%%%%%%%%%%%
	%Economista con vocación de servicio y compromiso con el desarrollo territorial equilibrado. Especializado en análisis socioeconómico regional, planificación del desarrollo y articulación de políticas multinivel (nacional, regional, local).
	
	%Experiencia en elaboración de planes de desarrollo concertado, diagnósticos territoriales, análisis de brechas de infraestructura y servicios, y diseño de programas de desarrollo productivo local. Capacidad de facilitar procesos participativos y articular actores públicos, privados y sociedad civil.
	
	%Busco oportunidades en gobiernos regionales, direcciones de planificación o agencias de desarrollo donde pueda contribuir a cerrar brechas de desarrollo, impulsar el crecimiento económico territorial sostenible y mejorar las condiciones de vida de las poblaciones más vulnerables.

	%%%%%%%%%%%%%%%%%%%%%%%%%%%%%%%%%%%%%%%%
	% VERSIÓN ALTERNATIVA - Gestión de información pública
	%%%%%%%%%%%%%%%%%%%%%%%%%%%%%%%%%%%%%%%%
	%Economista con especialización en gestión de información estadística y sistemas de información para el sector público. Experto en procesamiento, validación y análisis de bases de datos administrativas, registros institucionales y encuestas del sector público.
	
	%Dominio de estándares de calidad de datos, diseño de sistemas de indicadores, tableros de control para gestión pública y elaboración de boletines estadísticos institucionales. Familiarizado con interoperabilidad de sistemas, gobierno de datos y normativa de protección de datos personales.
	
	%Busco contribuir en áreas de estadística, planificación o tecnologías de información del sector público donde pueda aplicar mis conocimientos técnicos para mejorar la calidad de información institucional, fortalecer sistemas de seguimiento y evaluación, y facilitar la toma de decisiones basada en evidencia.

	%%%%%%%%%%%%%%%%%%%%%%%%%%%%%%%%%%%%%%%%
	% VERSIÓN ALTERNATIVA - Economía social
	%%%%%%%%%%%%%%%%%%%%%%%%%%%%%%%%%%%%%%%%
	%Economista con sensibilidad social y compromiso con la reducción de pobreza e inequidad. Especializado en programas sociales, economía del desarrollo y análisis de condiciones de vida de poblaciones vulnerables.
	
	%Experiencia en análisis de encuestas de hogares (ENAHO, ENAPRES), caracterización de pobreza multidimensional, evaluación de programas de transferencias condicionadas y estudios de inclusión social. Capacidad de traducir hallazgos técnicos en recomendaciones de política social accesibles.
	
	%Busco oportunidades en áreas de desarrollo social, inclusión económica o programas sociales donde pueda contribuir con mi formación técnica y mi compromiso ético para diseñar, implementar y evaluar intervenciones que mejoren las oportunidades y el bienestar de las poblaciones en situación de vulnerabilidad.

	%%%%%%%%%%%%%%%%%%%%%%%%%%%%%%%%%%%%%%%%
	% VERSIÓN ALTERNATIVA - Modernización del Estado
	%%%%%%%%%%%%%%%%%%%%%%%%%%%%%%%%%%%%%%%%
	%Economista con orientación en modernización de la gestión pública y reforma del Estado. Especializado en análisis de procesos institucionales, diseño de mejoras en la prestación de servicios públicos y simplificación administrativa.
	
	%Familiarizado con metodologías de gestión por procesos, mejora continua, gestión del cambio organizacional y transformación digital del Estado. Capacidad de elaborar diagnósticos institucionales, propuestas de rediseño organizacional y estrategias de fortalecimiento de capacidades.
	
	%Busco contribuir en áreas de modernización, desarrollo organizacional o gestión de la calidad del sector público donde pueda aplicar mi visión analítica y mi orientación a resultados para impulsar la eficiencia, transparencia y orientación ciudadana de las instituciones del Estado.
	
}
